Une entreprise qui fabrique des jouets doit effectuer des contrôles de conformité avant leur commercialisation. Dans cet exercice, on s'intéresse à deux tests effectués par l'entreprise de jouets : un test \emph{de fabrication} et un test \emph{de sécurité}.

\medskip

À la suite d'un grand nombre de vérifications, l'entreprise affirme que:

\begin{itemize}[label=$\bullet~$]
\item 95\,\% des jouets réussissent le test de fabrication ;
\item Parmi les jouets qui réussissent le test de fabrication, 98\,\% réussissent le test de sécurité ;
\item 1\,\% des jouets ne réussissent aucun des deux tests.
\end{itemize}

On choisit au hasard un jouet parmi les jouets produits. On note:

\begin{itemize}[label=$\bullet~$]
\item $F$ l'évènement : \og le jouet réussit le test de fabrication \fg{} ;
\item $S$ l'évènement : \og le jouet réussit le test de sécurité \fg.
\end{itemize}

\subsection*{Partie A}

\begin{enumerate}
\item À partir des données de l'énoncé, donner les probabilités $P(F)$ et $P_F(S)$.
\item
	\begin{enumerate}
		\item Construire un arbre pondéré qui illustre la situation avec les données disponibles dans l'énoncé.
		\item Montrer que $P_{\overline{F}}\left(\overline{S}\right) = 0,2$.
\end{enumerate}
\item Calculer la probabilité que le jouet choisi réussisse les deux tests.
\item Montrer que la probabilité que le jouet réussisse le test de sécurité vaut 0,97 arrondi au centième.
\item Lorsque le jouet a réussi le test de sécurité, quelle est la probabilité qu'il réussisse le test de fabrication? Donner une valeur approchée du résultat au centième.
\end{enumerate}

\subsection*{Partie B}

On prélève au hasard dans la production de l'entreprise un lot de $n$ jouets, où $n$ est un entier strictement positif. On suppose que ce prélèvement se fait sur une quantité suffisamment grande de jouets pour être assimilé à une succession de $n$ tirages indépendants avec remise.

On rappelle que la probabilité qu'un jouet réussisse le test de fabrication est égale à $0,95$.

Soit $S_n$ la variable aléatoire qui compte le nombre de jouets ayant réussi le test de fabrication. On
admet que $S_n$ suit la loi binomiale de paramètres $n$ et $p = 0,95$.

\medskip

\begin{enumerate}
\item Exprimer l'espérance et la variance de la variable aléatoire $S_n$ en fonction de $n$.
\item Dans cette question, on pose $n = 150$.
	\begin{enumerate}
		\item Déterminer une valeur approchée à $10^{-3}$ près de $P\left(S_{150} = 145\right)$. Interpréter ce résultat dans le contexte de l'exercice.
		\item Déterminer la probabilité qu'au moins 94\,\% des jouets de ce lot réussissent le test de
fabrication. Donner une valeur approchée du résultat à $10^{-3}$ près.
	\end{enumerate}
\item Dans cette question, l'entier naturel non nul $n$ n'est plus fixé.

Soit $F_n$ la variable aléatoire définie par : $F_n = \dfrac{S_n}{n}$. La variable aléatoire $F_n$ représente la proportion des jouets qui réussissent le test de fabrication dans un lot de $n$ jouets prélevés.

On note $E(F_n)$ l'espérance et $V(F_n)$ la variance de la variable aléatoire $F_n$.
	\begin{enumerate}
		\item Montrer que $E(F_n) = 0,95$ et que $V(F_n) = \dfrac{\np{0,0475}}{n}$.
		\item On s'intéresse à l'évènement $I$ suivant: \og la proportion de jouets qui réussissent le test de fabrication dans un lot de $n$ jouets est strictement comprise entre 93\,\% et 97\,\% \fg.

En utilisant l'inégalité de Bienaymé-Tchebychev, déterminer une valeur $n$ de la taille du lot
de jouets à prélever, à partir de laquelle la probabilité de l'évènement $I$ est supérieure ou égale à $0,96$.
	\end{enumerate}
\end{enumerate}
